\begin{table}[h]
	\centering
	\begin{tabular}{| c | c | c | c |}
 	\hline
Parameter(s)	& Prior		& Range    \\
\hline
$\sin^{2}\theta_{23}$	&  uniform 	& $[0.3; 0.7]$      \\
$\sin^{2} 2 \theta_{13}$ reactors		&  gaussian 		& $0.0857 \pm 0.0046$ 	\\
$\sin^{2} 2 \theta_{13}$ T2K only	   	&  uniform		& $[0.03; 0.2]$ 	\\
$\sin^{2} 2 \theta_{12}$  								&  gaussian 		& $0.846 \pm 0.021$ \\
$|$\dmsqtwothree$|$ (NH) / $|$\dmsqonethree$|$ (IH)      		&  uniform		& $ [2.3; 2.7] \times 10^{-3}~\mbox{eV}^2/\mbox{c}^4$		 \\
$\Delta m^{2}_{21}$      								&  gaussian 		& $(7.53 \pm 0.18) \times 10^{-5}~\mbox{eV}^2/\mbox{c}^4$		\\
$\delta_{CP}$                                               				&  uniform 	& $[-\pi; +\pi]$ \\
Mass Hierarchy										&  fixed 		& NH or IH \\  
\hline
\end{tabular}
\caption{
Treatment of the oscillation parameters in the $\nu/\bar{\nu}$ joint analysis. All the gaussian priors are from \cite{pdg_2014}. The parameter $\sin^{2} 2 \theta_{13}$ can be constrained with the measurement of reactor experiments using the gaussian prior, otherwise a uniform prior is used and the measurement relies only on the T2K data. The nominal values of $\sin^2 2 \theta_{13}$, $\sin^2 2 \theta_{12}$ and $\Delta m^2_{21}$ are from \cite{pdg_2014}, while all the other oscillation parameter values corresponds to the most probable values obtained by the Bayesian analysis on the T2K run 1-4 neutrino mode data \cite{T2KLongPaper}. The mass hierarchy is not marginalised but fixed to either NH or IH.
}
\label{tab:osc_param_priors}
\end{table}
